% ===========================================================================
\subsection{\file{new-design-reference/netlify.toml}}
% ===========================================================================

\textbf{TOML, 11 lines.} Instructions for a hosting service called Netlify: how to
publish the folder, and four security headers to send with every response.

\begin{provenance}{new-design-reference/netlify.toml}
Hand-written, added in the first commit \code{1f7ed12} and never edited. It is
\textbf{byte for byte identical} to the sibling repository's copy, which was
verified by comparing the two files directly: \code{diff} reports no differences at
all. Every other file in this pair differs somewhere; this one does not.

That is itself the evidence for how it got here: it is the one file in the design
that carried no site-specific content, so it was reused without needing a single
edit. The structure follows Netlify's own documented configuration format.
\textbf{Confidence: certain for hand-written and for the identity with the sibling,
both checked directly; no specific template is claimed.}
\end{provenance}

\subsubsection*{Why it exists}

Today, for no reason at all: nothing reads it.

The honest reading of the history is that the site was first expected to be hosted
on Netlify, a common choice for a static page like this, and the plan changed two
commits later. Commit \code{692e162} on 2026-07-28 is titled in part ``VPS rsync
deploy'', which is the moment Netlify stopped being the destination. This file was
never removed.

\begin{jargon}{TOML}
A file format for configuration designed to be obvious to read. A name in square
brackets starts a section; \code{name = "value"} sets a setting. Double square
brackets, as on line 5, mean ``another entry in a list of these'', so a file can
carry several header rules.
\end{jargon}

\subsubsection*{Line by line}

\begin{lstlisting}[firstnumber=1,caption={\file{netlify.toml}, lines 1--11}]
[build]
  publish = "."
  command = ""

[[headers]]
  for = "/*"
  [headers.values]
    X-Frame-Options = "DENY"
    X-Content-Type-Options = "nosniff"
    Referrer-Policy = "strict-origin-when-cross-origin"
    Permissions-Policy = "camera=(), microphone=(), geolocation=()"
\end{lstlisting}

Lines 1 to 3 describe how to build the site. \code{publish = "."} means publish this
directory itself, and \code{command = ""} means there is no build step to run first.
Those two lines are the entire ``no build step'' decision, stated to a hosting
provider.

Line 5 opens a headers rule and line 6 applies it to \code{/*}, every path on the
site.

\begin{jargon}{HTTP header}
An extra line of information a web server sends alongside a file, which the visitor
never sees. Some describe the file; the four below are instructions to the browser
about what it is allowed to do with the page.
\end{jargon}

Line 8, \code{X-Frame-Options = "DENY"}, forbids any other site from embedding this
page inside a frame. That defends against clickjacking, where an attacker loads your
real page invisibly over their own and tricks a visitor into clicking your buttons
while believing they are clicking something else.

Line 9, \code{X-Content-Type-Options = "nosniff"}, tells the browser to believe the
declared type of every file rather than guessing from its contents. Guessing is a
historical browser behaviour that can be manipulated into treating an uploaded image
as a script.

Line 10, \code{Referrer-Policy = "strict-origin-when-cross-origin"}, limits what is
disclosed when a visitor follows a link away. Within this site the full address is
passed; to another site only the domain; to an insecure site nothing at all. This is
the modern default and a sensible privacy choice.

Line 11, \code{Permissions-Policy}, switches off three device capabilities for this
page entirely: camera, microphone and location. The empty parentheses mean ``nobody,
not even this page''. The page is static text and gradients and needs none of them,
so turning them off means a future compromise cannot quietly request them.

\begin{keyidea}{These four headers are good practice and they are not in effect}
Every one is a sensible defensive choice, and the file they live in is read by a
service this project does not use. The page is served by Caddy on the agency VPS,
which never sees this file.

Moving these four lines into the Caddy configuration on the server would take a few
minutes and make them real. As things stand, the repository looks like it has
thought about security headers and the live site has none of them.
\end{keyidea}

\subsubsection*{What it connects to}

Nothing in this repository refers to it. Netlify appears nowhere else in the tree
and no workflow reads it. The deploy workflow does not even exclude it, so if that
workflow were restored this file would be copied to the VPS where Caddy would
ignore it.

\subsubsection*{Concepts introduced}

TOML as a configuration format, and the four HTTP response headers above. No
concept here appears anywhere else in the repository, which is another way of
saying this file is disconnected from everything.

\subsubsection*{How it breaks}

It cannot break, because nothing reads it. The realistic failure is the opposite: a
future maintainer reads this file, concludes the site sends these headers, and
builds on an assumption that is false.

\begin{honest}{An unused configuration file is worse than no configuration file}
This is a small file and a real trap. It states four security guarantees the live
site does not provide, and nothing in the repository indicates it is inert. Either
delete it, or move its contents to the Caddy configuration and delete it. Leaving it
is the only option with no upside.

The identical file sits in the sibling repository with the identical problem, so
fixing one should mean fixing both.
\end{honest}
